% TOP_PROBLEM: {"id": 4700008, "title": "A new Hilbert sixteenth-type problem", "category_id": 11, "category": {"id": 11, "name": "dynamical_systems", "display_name": "Dynamical Systems", "description": "Problems about long-term behavior of deterministic systems, Hamiltonian dynamics, and periodic orbits.", "slug": "dynamical-systems", "order_index": 11, "created_at": "2026-07-31T15:26:25.671Z"}, "status": "open", "problem_number": "AMR-046-0008", "set_id": 13, "statusQualification": "The least-monomial subquestion is solved with value4; the general bounds question remains the target."}
% FIRST_POSTED: 2026-10-02
% ENTRY_KIND: statement_only
% UnsolvedMath ID 4700008; source catalogue code AMR-046-0008
% !TeX program = lualatex
% Definition and sources only; this file is not an LLM solution attempt.
\documentclass[11pt,a4paper]{article}
\usepackage[margin=25mm]{geometry}
\usepackage{fontspec}
\setmainfont{lmroman10-regular.otf}[BoldFont=lmroman10-bold.otf,
 ItalicFont=lmroman10-italic.otf,BoldItalicFont=lmroman10-bolditalic.otf]
\usepackage{amsmath,amssymb,mathtools}
\usepackage{unicode-math}
\setmathfont{latinmodern-math.otf}
\AtBeginDocument{\let\setminus\smallsetminus}
\usepackage{xurl,hyperref}
\hypersetup{colorlinks=true,urlcolor=blue}
\setcounter{secnumdepth}{0}
\setlength{\parindent}{0pt}
\setlength{\parskip}{5pt}
\setlength{\emergencystretch}{3em}
\begin{document}
\section*{A new Hilbert sixteenth-type problem}\label{4700008}
{\small UnsolvedMath ID 4700008 · AMR-046-0008 · Dynamical Systems · AMR Open Problem Lists}
\subsection{Definitions and mathematical statement}
A monomial vector field on the plane is either $(x^ay^b,0)$ or
$(0,x^ay^b)$, where $a,b$ are nonnegative integers. A nonzero scalar
coefficient times each such term is counted once; identical powers in
different components are distinct terms. Let $\mathcal M_m$ consist of
polynomial vector fields obtained using at most $m$ distinct terms, with no
restriction on their degrees. Put
\[
 H^M(m)=\sup_{X\in\mathcal M_m}\#\{\text{limit cycles of }X\}.
\]
A limit cycle is an isolated nonconstant periodic orbit, counted once
regardless of its parametrization or multiplicity. The value $+\infty$ is
allowed in this supremum.

The main problem is to find useful upper and lower bounds on $H^M(m)$ and,
where possible, determine its exact values or growth. The original source
also asks for the least number of monomials allowing at least $m+1$ cycles.
That latter question is now resolved: later work gives no cycles for
$m\le3$ and at least twelve for $m=4$, so its least value is four.

The broader bounds question remains the target of this definition. In
particular a degree-dependent estimate does not answer it, because arbitrarily
large exponents are permitted while the number of terms stays fixed. The
2024 quadratic-growth lower bound is progress on this question, rather than
an upper bound or an exact determination of $H^M(m)$.
\subsection{Short English statement}
Bound the number of planar limit cycles by the number of monomial terms, with arbitrary exponents.
\subsection{Sources}
\begin{itemize}
\item[S1] ULAM.AI and the UnsolvedMath Contributors, supplied problems.json, record 4700008 (AMR-046-0008), AMR Open Problem Lists. Catalogue identity and source transcription; CC BY 4.0.
\url{https://huggingface.co/datasets/ulamai/UnsolvedMath}.
\item[S2] Gasull - Some open problems in low dimensional dynamical systems (2020), Problem environment 8: A new Hilbert sixteenth-type problem. Original source indexed by AMR-046-0008.
\url{https://arxiv.org/abs/2012.02524}.
\item[S3] Gasull and Santana, On a variant of Hilbert’s 16th problem (2024), Theorems 1–2; partial progress and the resolved least-term subquestion.
\url{https://arxiv.org/abs/2405.04281}.
\end{itemize}
\end{document}
